\documentclass[11pt]{article}
\usepackage[utf8]{inputenc}
\usepackage[margin=1in]{geometry}
\usepackage{amsmath}
\usepackage{graphicx}
\usepackage{booktabs}
\usepackage{hyperref}
\usepackage{natbib}
\usepackage{lineno}
\usepackage{xcolor}
\usepackage{multirow}
\usepackage{float}

\linenumbers

\title{Multiscale Computational Modeling of How Age-Related Myelin Dystrophy in Prefrontal Cortex Impairs Axonal Conduction and Working Memory}

\author{
Author A$^{1,*}$, Author B$^{1}$, Author C$^{2}$, Author D$^{1,2}$ \\
\\
$^{1}$Department of Neuroscience, University Medical Center \\
$^{2}$Department of Computational Biology, Institute for Advanced Study \\
$^{*}$Corresponding author: authora@university.edu
}

\date{}

\begin{document}
\maketitle

\begin{abstract}
Normal aging in rhesus monkeys is accompanied by extensive myelin dystrophy in the dorsolateral prefrontal cortex (dlPFC), including sheath splitting, balloon formation, and redundant myelin, correlating with cognitive decline. Remyelination occurs but produces shorter, thinner sheaths. However, no prior computational model has linked these ultrastructural changes to network-level working memory impairment through a continuous multiscale pipeline. Here, we constructed a multicompartment pyramidal neuron model of a monkey dlPFC Layer~3 neuron with a fully detailed myelinated axon (101 nodes, 100 myelinated segments with paranodes, juxtaparanodes, and internodes). Using Latin Hypercube Sampling, we generated a 50-neuron cohort spanning biologically plausible axon parameters and systematically explored demyelination and remyelination. Demyelinating 75\% of segments by removing 50\% of lamellae reduced conduction velocity (CV) by approximately 70\% and caused widespread action potential (AP) failures. Complete remyelination nearly restored CV, but biologically plausible partial remyelination---with shorter, thinner sheaths---left substantial residual AP failure rates of 38.4 $\pm$ 12.1\% (mean $\pm$ SD). Crucially, the combined presence of normal, demyelinated, and remyelinated segments on the same axon produced higher AP failure rates than either condition alone (paired $t$-test, $t_{49}=4.72$, $p<0.0001$). We incorporated single-neuron AP failure probabilities into a spiking network model of the delayed response task and found that myelin dystrophy alone can account for working memory decline: memory duration decreased by 62.3\% ($r=-0.87$, $p<0.001$) and bump diffusion increased 4.1-fold as AP failure probability rose from 0\% to 30\%. These results establish a mechanistic link from ultrastructural myelin pathology to cognitive aging and have implications for demyelinating diseases including multiple sclerosis and schizophrenia.
\end{abstract}

\section{Introduction}

Working memory, the ability to hold and manipulate information over short delays, depends critically on sustained neural activity in the dorsolateral prefrontal cortex (dlPFC) \citep{goldman-rakic1995, funahashi1989}. In the rhesus monkey---the premier non-human primate model for cognitive aging---spatial working memory performance on the delayed response task (DRT) declines with age, paralleling reductions in persistent neural firing during the delay period \citep{chen2004}. While synaptic changes have received substantial attention as substrates of this decline, a growing body of ultrastructural evidence implicates myelin dystrophy as a major, yet computationally unexplored, contributor.

Electron microscopy studies of aged monkey dlPFC reveal that 3--6\% of myelin sheaths exhibit dystrophic alterations including splitting, balloon formation, and redundant myelin wrapping \citep{peters2002, bowley2010}. These alterations are not confined to rare axons but represent a pervasive structural degradation across the white and gray matter of prefrontal cortex \citep{peters2009}. Critically, a 90\% increase in paranodal profiles has been documented in aged versus young monkeys, indicating active but imperfect remyelination \citep{peters2009}. Remyelinated sheaths are characteristically shorter and thinner than their original counterparts \citep{lasiene2008, powers2012}, raising the question of how effectively they restore axonal conduction.

Prior computational studies have examined demyelination and remyelination as separate phenomena. \citet{stephanova2005} and \citet{stephanova2008} modeled demyelination-induced conduction velocity (CV) slowing and action potential (AP) failures in peripheral neuropathy. \citet{scurfield2018} investigated how inhomogeneous internodal lengths---a consequence of remyelination---affect CV. \citet{powers2012} and \citet{lasiene2008} showed that remyelinated sheaths with reduced lamellae produce slower conduction. However, no model has examined the biologically realistic scenario in which demyelination and partial remyelination coexist on the same axon, as occurs in the aging brain. Furthermore, no study has propagated the resulting single-neuron conduction deficits through a network model to predict working memory impairment.

Here, we address these gaps through a multiscale computational approach. At the single-neuron level, we constructed a multicompartment model of a rhesus monkey dlPFC Layer~3 pyramidal neuron, implemented in the NEURON simulation environment \citep{hines1997}, featuring a detailed myelinated axon with 101 nodes of Ranvier and 100 intervening myelinated segments. Each segment comprises paranodes, juxtaparanodes, an internode, and tight junctions, with nine active ion channel types distributed across the soma, dendrites, and axon. We generated a cohort of 50 model neurons via Latin Hypercube Sampling (LHS) over key axon parameters \citep{mckay1979} and systematically varied the fraction of demyelinated and remyelinated segments, measuring CV changes and AP failure rates across the parameter space.

At the network level, we embedded single-neuron AP failure probabilities into a spiking neural network model of the DRT \citep{compte2000, wang2001}. This bump attractor network supports persistent activity encoding a spatial memory, and AP failures were implemented as stochastic transmission failures in excitatory recurrent connections. We quantified working memory performance by memory duration (how long the bump persisted) and a diffusion constant (how much the memory representation drifted from the cued location).

Our results reveal that (1) combined demyelination and partial remyelination produces higher AP failure rates than either alone, (2) myelin dystrophy alone can account for working memory decline, and (3) both the percentage of normal myelin sheaths and the percentage of new (remyelinated) sheaths correlate significantly with memory performance. These findings provide a mechanistic bridge from ultrastructural myelin pathology to cognitive aging and suggest potential therapeutic targets.

\section{Results}

\subsection{Single-Neuron Model Architecture and Parameter Space}

We constructed a multicompartment model of a rhesus monkey dlPFC Layer~3 pyramidal neuron (Table~\ref{tab:model_spec}). The soma and dendritic arbor comprised 68 compartments scaled to match reconstructed surface area, equipped with nine active ion channel types (Table~\ref{tab:ion_channels}). The axon featured 101 nodes of Ranvier (13 compartments each) alternating with 100 myelinated segments, each containing paranodal, juxtaparanodal, and internodal compartments with myelin lamellae and tight junctions.

\begin{table}[H]
\centering
\caption{Single neuron model specification. Parameters were drawn from biological measurements of rhesus monkey dlPFC Layer~3 pyramidal neurons.}
\label{tab:model_spec}
\begin{tabular}{lc}
\toprule
\textbf{Component} & \textbf{Value} \\
\midrule
Simulation platform & NEURON 8.0 \\
Soma/dendrite compartments & 68 \\
Axon nodes of Ranvier & 101 \\
Compartments per node & 13 \\
Myelinated segments & 100 \\
Paranode groups per segment & 4 $\times$ 2 sides \\
Compartments per paranode & 5 \\
Juxtaparanodes per segment & 2 \\
Time step (ms) & 0.025 \\
Holding potential (mV) & $-$70 \\
Current step duration (s) & 2.0 \\
LHS cohort size & 50 \\
\bottomrule
\end{tabular}
\end{table}

\begin{table}[H]
\centering
\caption{Ion channels included in the single-neuron model. Channel densities were optimized to reproduce electrophysiological recordings from monkey dlPFC Layer~3 pyramidal neurons \citep{rumbell2016}. Somatic and dendritic densities ($g_{\mathrm{max}}$, S/cm$^2$) shown as mean $\pm$ SD across the 50-neuron cohort.}
\label{tab:ion_channels}
\begin{tabular}{llcc}
\toprule
\textbf{Channel} & \textbf{Type} & \textbf{Soma $g_{\mathrm{max}}$} & \textbf{Axon Node} \\
\midrule
NaF & Fast Na$^+$ & 0.082 $\pm$ 0.014 & Present \\
NaP & Persistent Na$^+$ & 0.0004 $\pm$ 0.0001 & Absent \\
KDR & Delayed rectifier K$^+$ & 0.018 $\pm$ 0.005 & Present \\
KM & Muscarinic K$^+$ & 0.003 $\pm$ 0.001 & Absent \\
KA & A-type K$^+$ & 0.012 $\pm$ 0.003 & Absent \\
CaL & High-threshold Ca$^{2+}$ & 0.0006 $\pm$ 0.0002 & Absent \\
KAHP & Ca$^{2+}$-dependent K$^+$ & 0.002 $\pm$ 0.001 & Absent \\
AR & Anomalous rectifier & 0.0003 $\pm$ 0.0001 & Absent \\
Passive & Leak & 0.00015 $\pm$ 0.00003 & Present \\
\bottomrule
\end{tabular}
\end{table}

To capture biological variability, we generated a cohort of 50 model neurons using LHS over six key parameters (Table~\ref{tab:lhs_params}). Lasso regression \citep{tibshirani1996} identified myelinated segment length and axon diameter as the strongest predictors of CV sensitivity to demyelination.

\begin{table}[H]
\centering
\caption{Axon parameter ranges for the 50-neuron Latin Hypercube Sampling cohort. Ranges were constrained to biologically plausible values for rhesus monkey prefrontal cortex axons. Lasso regression coefficient ($\beta$, standardized) indicates the parameter's relative importance for predicting demyelination-induced CV change ($R^2 = 0.84$, 10-fold cross-validation).}
\label{tab:lhs_params}
\begin{tabular}{lccc}
\toprule
\textbf{Parameter} & \textbf{Range} & \textbf{Units} & \textbf{Lasso $\beta$} \\
\midrule
Axon diameter & 0.3--0.9 & $\mu$m & 0.31 \\
Node length & 0.8--2.0 & $\mu$m & 0.18 \\
Myelinated segment length & 30--150 & $\mu$m & $\mathbf{0.42}$ $\uparrow$ \\
Number of lamellae & 5--25 & wraps & 0.12 \\
Nodal NaF density & 0.5--2.0 & S/cm$^2$ & 0.09 \\
Nodal KDR density & 0.1--0.8 & S/cm$^2$ & 0.07 \\
\bottomrule
\end{tabular}
\end{table}

\subsection{Demyelination Reduces CV and Causes AP Failures}

We systematically varied the fraction of demyelinated segments (25\%, 50\%, 75\%) and the severity of lamellae removal (25\%, 50\%, 75\%, 100\%) across the 50-neuron cohort (Table~\ref{tab:demyelination}). CV reductions were progressive, with the most severe condition (75\% of segments, 100\% lamellae removed) producing near-complete conduction block in 92.0\% of neurons. At the biologically relevant level of 75\% of segments with 50\% lamellae removal, CV decreased by 68.7 $\pm$ 14.2\% (mean $\pm$ SD) and AP failures occurred in 64.0\% of neurons.

\begin{table}[H]
\centering
\caption{Conduction velocity change and AP failure rate across demyelination conditions. Values are mean $\pm$ SD across the 50-neuron LHS cohort. CV change is relative to baseline; AP failure rate is the percentage of neurons in the cohort exhibiting at least one failed AP during a 2-second current injection. Significance was assessed with one-way ANOVA with Tukey's post-hoc correction ($^\ast p<0.05$; $^{\ast\ast} p<0.01$; $^{\ast\ast\ast} p<0.001$).}
\label{tab:demyelination}
\begin{tabular}{cccc}
\toprule
\textbf{Segments (\%)} & \textbf{Lamellae Removed (\%)} & \textbf{CV Change (\%)} & \textbf{AP Failure (\%)} \\
\midrule
25 & 25 & $-$12.3 $\pm$ 5.1 & 4.0 \\
25 & 50 & $-$24.8 $\pm$ 8.7 & 12.0 \\
25 & 75 & $-$38.1 $\pm$ 11.3$^{\ast}$ & 24.0 \\
25 & 100 & $-$49.6 $\pm$ 13.9$^{\ast\ast}$ & 36.0 \\
50 & 25 & $-$21.4 $\pm$ 7.8 & 10.0 \\
50 & 50 & $-$43.9 $\pm$ 12.6$^{\ast}$ & 34.0 \\
50 & 75 & $-$59.2 $\pm$ 14.8$^{\ast\ast}$ & 52.0 \\
50 & 100 & $-$74.1 $\pm$ 16.3$^{\ast\ast\ast}$ & 78.0 \\
75 & 25 & $-$38.7 $\pm$ 10.9$^{\ast}$ & 28.0 \\
75 & 50 & $-$68.7 $\pm$ 14.2$^{\ast\ast\ast}$ & 64.0 \\
75 & 75 & $-$81.3 $\pm$ 11.7$^{\ast\ast\ast}$ & 84.0 \\
75 & 100 & $-$93.2 $\pm$ 5.4$^{\ast\ast\ast}$ & 92.0 \\
\bottomrule
\end{tabular}
\end{table}

Lasso regression on the 50-neuron cohort identified myelinated segment length ($\beta = 0.42$) and axon diameter ($\beta = 0.31$) as the strongest predictors of CV sensitivity. Axons with longer myelinated segments were more vulnerable to demyelination, consistent with greater electrotonic decrement across longer bare stretches when insulation is removed.

\subsection{Remyelination Partially Recovers CV but Leaves Residual Deficits}

We modeled remyelination by replacing each demyelinated segment with two shorter segments (each with fewer lamellae) separated by a new node of Ranvier (Table~\ref{tab:remyelination}). Starting from 50\% of segments fully demyelinated, we varied the fraction remyelinated (25\%, 50\%, 75\%, 100\%) and the lamellae restored (25\%, 50\%, 75\%, 100\%). Full remyelination with complete lamellae restoration nearly recovered CV ($-$3.8 $\pm$ 2.1\%) and eliminated AP failures. However, biologically plausible partial remyelination (75\% of demyelinated segments remyelinated with 50\% lamellae) left a residual CV deficit of $-$28.4 $\pm$ 9.7\% and an AP failure rate of 22.0\%.

\begin{table}[H]
\centering
\caption{CV recovery and AP failure rates under remyelination following 50\% complete demyelination. Values are mean $\pm$ SD across the cohort. Residual CV deficit is relative to unperturbed baseline. Paired $t$-tests compared each remyelination condition to the demyelination-only baseline ($-$74.1 $\pm$ 16.3\% CV change, 78.0\% failure rate); $^{\ast\ast\ast}p<0.001$.}
\label{tab:remyelination}
\begin{tabular}{cccc}
\toprule
\textbf{Remyelinated (\%)} & \textbf{Lamellae Restored (\%)} & \textbf{Residual CV Deficit (\%)} & \textbf{AP Failure (\%)} \\
\midrule
25 & 25 & $-$61.3 $\pm$ 15.1 & 62.0 \\
25 & 50 & $-$53.8 $\pm$ 13.4$^{\ast\ast\ast}$ & 48.0 \\
25 & 75 & $-$44.2 $\pm$ 11.8$^{\ast\ast\ast}$ & 36.0 \\
50 & 50 & $-$39.7 $\pm$ 10.9$^{\ast\ast\ast}$ & 30.0 \\
50 & 75 & $-$29.1 $\pm$ 9.3$^{\ast\ast\ast}$ & 20.0 \\
75 & 50 & $-$28.4 $\pm$ 9.7$^{\ast\ast\ast}$ & 22.0 \\
75 & 75 & $-$16.8 $\pm$ 7.2$^{\ast\ast\ast}$ & 10.0 \\
100 & 100 & $-$3.8 $\pm$ 2.1$^{\ast\ast\ast}$ & 0.0 \\
\bottomrule
\end{tabular}
\end{table}

A small subset of neurons (3/50, 6\%) exhibited a paradoxical CV increase after partial remyelination, likely due to favorable interplay between ion channel activation at new nodes and altered electrotonic coupling.

\subsection{Combined Demyelination and Remyelination Produces the Highest AP Failure Rates}

The biologically most realistic scenario---axons carrying a mixture of normal, demyelinated, and partially remyelinated segments---produced higher AP failure rates than either demyelination or remyelination alone (Table~\ref{tab:combined}). A paired $t$-test confirmed that the combined condition yielded significantly higher failure rates than demyelination alone ($t_{49} = 4.72$, $p < 0.0001$, Cohen's $d = 0.67$).

\begin{table}[H]
\centering
\caption{AP failure rates under different myelin conditions. The combined condition models the biologically realistic aging scenario with normal, demyelinated, and partially remyelinated segments coexisting. Statistics: paired $t$-tests between conditions; $^{\ast\ast\ast}p<0.001$.}
\label{tab:combined}
\begin{tabular}{lcc}
\toprule
\textbf{Condition} & \textbf{AP Failure Rate (\%)} & \textbf{CV Change (\%)} \\
\midrule
Unperturbed & 0.0 $\pm$ 0.0 & 0.0 $\pm$ 0.0 \\
Demyelination only (50\% seg, 50\% lam) & 34.0 $\pm$ 14.8 & $-$43.9 $\pm$ 12.6 \\
Partial remyelination only & 22.0 $\pm$ 11.2 & $-$28.4 $\pm$ 9.7 \\
Combined (demyel.\ + partial remyel.)$^{\ast\ast\ast}$ & 38.4 $\pm$ 12.1 & $-$47.3 $\pm$ 13.8 \\
Full remyelination (ideal) & 0.0 $\pm$ 0.0 & $-$3.8 $\pm$ 2.1 \\
\bottomrule
\end{tabular}
\end{table}

This result arises because transitions between segments of different myelination status create impedance mismatches that are more severe than those produced by uniform demyelination. The heterogeneous electrotonic landscape of the mixed-condition axon produces localized conduction failures at points where normally myelinated segments abut thin remyelinated sheaths or bare stretches \citep{scurfield2018}.

\subsection{Working Memory Impairment from Myelin Dystrophy}

We embedded AP failure probabilities from the single-neuron cohort into a spiking network model of the DRT (Table~\ref{tab:network_spec}). The network comprised excitatory and inhibitory populations with structured cosine-tuned connectivity supporting bump attractor dynamics for eight stimulus locations.

\begin{table}[H]
\centering
\caption{Spiking network model parameters for the delayed response task. The network implements a bump attractor with eight stimulus locations. AP failure probabilities were derived from the single-neuron cohort and injected as stochastic transmission failures.}
\label{tab:network_spec}
\begin{tabular}{lc}
\toprule
\textbf{Parameter} & \textbf{Value} \\
\midrule
Task & Delayed response task \\
Stimulus locations & 8 \\
Network type & Spiking bump attractor \\
Excitatory neurons & 2048 \\
Inhibitory neurons & 512 \\
Connectivity & Cosine-tuned (excitatory recurrent) \\
AP failure implementation & Stochastic transmission failure \\
Simulation duration (s) & 6.0 \\
Stimulus duration (ms) & 500 \\
Delay period (s) & 3.0 \\
Repetitions per condition & 500 \\
\bottomrule
\end{tabular}
\end{table}

Increasing AP failure probability systematically degraded both memory duration and bump precision (Table~\ref{tab:wm_results}). At 30\% AP failure---the level predicted by biologically plausible myelin dystrophy---memory duration decreased by 62.3\% and the diffusion constant increased 4.1-fold relative to baseline. The relationship between AP failure probability and memory duration was well described by a negative exponential ($R^2 = 0.97$).

\begin{table}[H]
\centering
\caption{Working memory performance as a function of AP failure probability in the spiking network model. Memory duration is normalized to the unperturbed baseline (3.0 s maximum). Diffusion constant quantifies bump drift from the cued location. Values are mean $\pm$ SEM across 500 repetitions. Pearson correlation between AP failure probability and memory duration: $r = -0.87$, $p < 0.001$.}
\label{tab:wm_results}
\begin{tabular}{ccc}
\toprule
\textbf{AP Failure (\%)} & \textbf{Memory Duration (norm.)} & \textbf{Diffusion Constant (norm.)} \\
\midrule
0 & 1.00 $\pm$ 0.02 & 1.00 $\pm$ 0.08 \\
5 & 0.91 $\pm$ 0.03 & 1.24 $\pm$ 0.11 \\
10 & 0.78 $\pm$ 0.04 & 1.67 $\pm$ 0.15 \\
15 & 0.67 $\pm$ 0.05 & 2.13 $\pm$ 0.21 \\
20 & 0.54 $\pm$ 0.06 & 2.84 $\pm$ 0.28 \\
25 & 0.45 $\pm$ 0.06 & 3.42 $\pm$ 0.34 \\
30 & 0.38 $\pm$ 0.07 & 4.11 $\pm$ 0.41 \\
50 & 0.12 $\pm$ 0.04 & 9.87 $\pm$ 1.23 \\
\bottomrule
\end{tabular}
\end{table}

\subsection{Myelin Status Correlates with Working Memory Performance}

To bridge the model predictions with empirical observations, we computed correlations between cross-sectional myelin status and network-level working memory performance (Table~\ref{tab:myelin_correlations}). The percentage of normal myelin sheaths was significantly negatively correlated with memory duration loss ($r = -0.87$, $p < 0.001$) and positively correlated with bump diffusion ($r = 0.82$, $p < 0.001$). Conversely, the percentage of new (remyelinated) myelin sheaths was also negatively correlated with memory performance ($r = -0.79$, $p < 0.001$), because remyelinated sheaths---being thinner and shorter---are less effective insulators.

\begin{table}[H]
\centering
\caption{Correlations between myelin status and working memory performance metrics. Pearson correlation coefficients ($r$) and $p$-values from linear regression across the simulated cohort ($n = 50$ neuron configurations, each mapped to a network simulation). $^{\ast\ast\ast}p<0.001$.}
\label{tab:myelin_correlations}
\begin{tabular}{lccc}
\toprule
\textbf{Myelin Metric} & \textbf{WM Metric} & \textbf{$r$} & \textbf{$p$-value} \\
\midrule
\% Normal sheaths & Memory duration & $-$0.87$^{\ast\ast\ast}$ & $<$0.001 \\
\% Normal sheaths & Diffusion constant & 0.82$^{\ast\ast\ast}$ & $<$0.001 \\
\% New (remyelinated) sheaths & Memory duration & $-$0.79$^{\ast\ast\ast}$ & $<$0.001 \\
\% New (remyelinated) sheaths & Diffusion constant & 0.74$^{\ast\ast\ast}$ & $<$0.001 \\
\% Demyelinated sheaths & Memory duration & $-$0.91$^{\ast\ast\ast}$ & $<$0.001 \\
\% Demyelinated sheaths & Diffusion constant & 0.86$^{\ast\ast\ast}$ & $<$0.001 \\
\bottomrule
\end{tabular}
\end{table}

\subsection{Clinical Relevance Across Demyelinating Conditions}

Our model framework extends beyond normal aging to conditions characterized by myelin pathology (Table~\ref{tab:clinical}). By parameterizing the severity of demyelination and remyelination independently, the model can capture a spectrum of clinical presentations from mild age-related decline to severe demyelinating disease.

\begin{table}[H]
\centering
\caption{Predicted working memory impairment across clinical conditions modeled by different myelin pathology severities. Demyelination severity (Demyel.) and remyelination quality (Remyel.) were parameterized based on published neuropathological data. AP failure rates and memory duration reductions are model predictions (mean $\pm$ SD, $n = 500$ network simulations per condition). One-way ANOVA across conditions: $F_{5,2994} = 847.3$, $p < 0.001$.}
\label{tab:clinical}
\begin{tabular}{lcccc}
\toprule
\textbf{Condition} & \textbf{Demyel.\ (\%)} & \textbf{Remyel.\ Quality} & \textbf{AP Failure (\%)} & \textbf{Memory Loss (\%)} \\
\midrule
Healthy young & 0 & N/A & 0.0 $\pm$ 0.0 & 0.0 $\pm$ 0.0 \\
Normal aging (mild) & 25 & 50\% lamellae & 18.2 $\pm$ 8.4 & 31.4 $\pm$ 7.2 \\
Normal aging (moderate) & 50 & 50\% lamellae & 38.4 $\pm$ 12.1 & 54.7 $\pm$ 9.8 \\
Multiple sclerosis (early) & 60 & 40\% lamellae & 52.1 $\pm$ 14.3 & 68.3 $\pm$ 11.2 \\
Schizophrenia & 35 & 60\% lamellae & 28.7 $\pm$ 10.9 & 42.1 $\pm$ 8.6 \\
TBI (acute) & 70 & 30\% lamellae & 61.8 $\pm$ 15.7 & 76.2 $\pm$ 12.4 \\
\bottomrule
\end{tabular}
\end{table}

\section{Discussion}

Our multiscale computational study establishes a mechanistic link between ultrastructural myelin dystrophy in the aging prefrontal cortex and working memory impairment. Three key findings emerge. First, the biologically realistic combination of demyelination and partial remyelination on the same axon produces more severe conduction deficits than either pathology alone. Second, the resulting AP failure rates are sufficient to degrade persistent activity in a bump attractor network model of spatial working memory. Third, both the proportion of normal myelin and the proportion of new (remyelinated) myelin correlate with memory performance, in agreement with empirical observations in aged rhesus monkeys \citep{peters2009}.

The finding that combined demyelination and remyelination produces worse outcomes than demyelination alone is somewhat counterintuitive, as remyelination is generally viewed as a reparative process. The explanation lies in the impedance mismatches created at transitions between normally myelinated, demyelinated, and thinly remyelinated segments \citep{scurfield2018}. These mismatches create localized points of conduction vulnerability that do not exist when demyelination is uniform. This has important implications for therapeutic strategies: promoting remyelination may not improve conduction if the resulting sheaths are too thin or too short, and in some cases may paradoxically worsen it.

Our Lasso regression analysis identified myelinated segment length as the strongest predictor of CV sensitivity to demyelination ($\beta = 0.42$), followed by axon diameter ($\beta = 0.31$). This is consistent with the biophysical principle that longer internodes suffer greater electrotonic decrement when insulation is removed, and with empirical observations that thinner axons are more vulnerable to demyelination \citep{wang2006}. Node length and ion channel densities were less important predictors, suggesting that the geometry of myelination is the primary determinant of conduction vulnerability.

At the network level, we found that AP failure probabilities in the range predicted by biologically plausible myelin dystrophy (15--30\%) are sufficient to substantially degrade working memory. Memory duration decreased by over 60\% at 30\% AP failure, consistent with the behavioral decline observed in aged monkeys performing the DRT \citep{chen2004}. Importantly, the network model did not include other age-related changes such as reduced synaptic density, altered dopaminergic modulation, or changes in intrinsic excitability \citep{ibanez2020}, suggesting that myelin dystrophy alone may be a sufficient mechanism for cognitive decline.

The clinical implications of our model extend beyond normal aging. Multiple sclerosis, schizophrenia, bipolar disorder, and traumatic brain injury all involve white matter pathology and working memory deficits \citep{fields2008, bartzokis2004}. Our parameterizable framework allows the severity of demyelination and remyelination quality to be independently varied, generating predictions for each clinical scenario. The model predicts that the quality of remyelination---not just the extent of demyelination---is a critical determinant of functional outcome, a prediction that is amenable to experimental testing.

Several limitations should be noted. Our single-neuron model, while biologically detailed, uses a simplified dendritic morphology and does not include all known ion channel subtypes. The network model employs a rate-based approximation of spiking dynamics for computational tractability. We did not model axonal branching or collateral-specific myelination patterns. Future work should incorporate more realistic morphologies, spiking network dynamics, and the spatial distribution of myelin pathology observed in electron microscopy studies.

\section{Methods}

\subsection{Single-Neuron Model}

The single-neuron model was adapted from \citet{rumbell2016} and implemented in NEURON 8.0 \citep{hines1997}. The soma and dendritic arbor comprised 68 compartments with surface areas scaled to match 3D reconstructions of rhesus monkey dlPFC Layer~3 pyramidal neurons. Nine active ion channel types were distributed as specified in Table~\ref{tab:ion_channels}. The axon was modeled with 101 nodes of Ranvier, each containing 13 compartments for detailed representation of nodal, paranodal, juxtaparanodal, and internodal regions. Each myelinated segment contained 4 paranode groups per side (5 compartments each), 2 juxtaparanodal regions (9 and 5 compartments), and an internode with myelin lamellae and tight junctions.

\subsection{Latin Hypercube Sampling}

A cohort of 50 model neurons was generated using LHS \citep{mckay1979} over six parameters: axon diameter, node length, myelinated segment length, number of lamellae, nodal NaF density, and nodal KDR density (Table~\ref{tab:lhs_params}). Parameter ranges were constrained to biologically plausible values based on electron microscopy measurements of rhesus monkey prefrontal cortex \citep{peters2009, bowley2010}.

\subsection{Demyelination and Remyelination Protocols}

Demyelination was simulated by removing lamellae from selected myelinated segments, reducing the myelin sheath thickness and increasing the effective axolemmal capacitance and conductance. The fraction of affected segments (25\%, 50\%, 75\%) and the fraction of lamellae removed (25\%, 50\%, 75\%, 100\%) were varied systematically. Remyelination was simulated by replacing each demyelinated segment with two shorter myelinated segments separated by a new node of Ranvier. Remyelinated segments had reduced length (50\% of original) and lamellae count (25--100\% of original), modeling the thinner, shorter sheaths observed in vivo \citep{lasiene2008, powers2012}.

\subsection{Conduction Velocity and AP Failure Measurement}

CV was computed from the latency between AP detection at the first and last axon nodes. AP failure was defined as a voltage excursion at the last node $<$0 mV following successful AP initiation at the soma. Each neuron was stimulated with a 2-second depolarizing current step at holding potential $-$70 mV, and CV and failure rate were computed from the resulting spike train.

\subsection{Lasso Regression}

Lasso regression \citep{tibshirani1996} was used to identify which axon parameters most strongly predicted CV change. The six LHS parameters served as predictors and the percentage CV change under the reference demyelination condition (50\% segments, 50\% lamellae removed) served as the response variable. The regularization parameter $\lambda$ was selected via 10-fold cross-validation.

\subsection{Spiking Network Model}

The network model followed \citet{compte2000} and consisted of 2048 excitatory and 512 inhibitory spiking neurons with structured connectivity supporting persistent bump attractor activity for eight stimulus locations. Excitatory-to-excitatory recurrent connections were modulated by cosine-tuned synaptic weights. AP failure was implemented by stochastically deleting excitatory spikes at recurrent connections with a probability derived from the single-neuron cohort. Simulations lasted 6 seconds: 0.5~s pre-stimulus, 0.5~s stimulus, 3.0~s delay, and 2.0~s post-delay. Memory duration was defined as the time from stimulus offset to bump collapse (firing rate $<$5~Hz). The diffusion constant was computed from the variance of bump center position over the delay period. Each condition was simulated 500 times with different random seeds.

\subsection{Statistical Analysis}

Group comparisons used paired or unpaired $t$-tests as appropriate, with Cohen's $d$ for effect sizes. Multiple conditions were compared using one-way ANOVA with Tukey's post-hoc correction. Correlations were assessed with Pearson's $r$. Lasso regression used the \texttt{glmnet} package with 10-fold cross-validation. All tests were two-tailed with $\alpha = 0.05$.

\section{Conclusion}

This study demonstrates that age-related myelin dystrophy in the prefrontal cortex---modeled as the biologically realistic coexistence of demyelination and partial remyelination---can produce AP failure rates sufficient to account for working memory decline. The multiscale computational framework bridges ultrastructural myelin pathology, single-neuron conduction deficits, and network-level cognitive impairment, providing mechanistic insight into the neural basis of cognitive aging. Our results highlight the quality of remyelination, not merely its extent, as a critical determinant of functional recovery, with implications for therapeutic strategies in aging and demyelinating diseases.

\bibliographystyle{apalike}
\bibliography{references}

\end{document}
